\documentclass[9pt,a4paper,twocolumn]{ctexart}

% --- Packages ---
\usepackage[top=1.2cm, bottom=1.2cm, left=1.2cm, right=1.2cm, columnsep=0.8cm]{geometry}
\usepackage{hyperref}
\usepackage{graphicx}
\usepackage{float}
\usepackage{longtable}
\usepackage{tabularx}
\usepackage{booktabs}
\usepackage{enumitem}
\usepackage{xcolor}
\usepackage{fancyhdr}
\usepackage{listings}
\usepackage{fontspec}
\usepackage{titlesec}
\usepackage{caption}

% --- Code block styling ---
\definecolor{codebg}{RGB}{245,245,245}
\definecolor{codeframe}{RGB}{220,220,220}
\definecolor{codekw}{RGB}{0,0,192}
\definecolor{codecomment}{RGB}{0,128,0}
\definecolor{codestring}{RGB}{163,21,21}
\definecolor{codenum}{RGB}{128,0,128}
\definecolor{codepunct}{RGB}{90,90,90}
\definecolor{badgegreen}{RGB}{34,139,34}
\definecolor{badgeblue}{RGB}{30,144,255}
\definecolor{badgeblack}{RGB}{50,50,50}

% --- Difficulty badges (rounded corners via tikz, pre-rendered in saveboxes) ---
\usepackage{tikz}
\newsavebox{\badgechubox}
\savebox{\badgechubox}{\tikz[baseline]{\node[fill=badgegreen,text=white,rounded corners=2.5pt,inner xsep=4pt,inner ysep=1.5pt,font=\sffamily\footnotesize\bfseries]{初级};}}
\newsavebox{\badgezhongbox}
\savebox{\badgezhongbox}{\tikz[baseline]{\node[fill=badgeblue,text=white,rounded corners=2.5pt,inner xsep=4pt,inner ysep=1.5pt,font=\sffamily\footnotesize\bfseries]{中级};}}
\newsavebox{\badgegaobox}
\savebox{\badgegaobox}{\tikz[baseline]{\node[fill=badgeblack,text=white,rounded corners=2.5pt,inner xsep=4pt,inner ysep=1.5pt,font=\sffamily\footnotesize\bfseries]{高级};}}
\newcommand{\lvchu}{\usebox{\badgechubox}}
\newcommand{\lvzhong}{\usebox{\badgezhongbox}}
\newcommand{\lvgao}{\usebox{\badgegaobox}}
\pdfstringdefDisableCommands{%
  \def\lvchu{[初级]}%
  \def\lvzhong{[中级]}%
  \def\lvgao{[高级]}%
}

\lstset{
  basicstyle=\ttfamily\scriptsize,
  keywordstyle=\color{codekw}\bfseries,
  commentstyle=\color{codecomment}\itshape,
  stringstyle=\color{codestring},
  backgroundcolor=\color{codebg},
  frame=single,
  rulecolor=\color{codeframe},
  breaklines=true,
  breakatwhitespace=false,
  breakindent=0pt,
  postbreak=\mbox{\textcolor{red}{$\hookrightarrow$}\space},
  numbers=left,
  numberstyle=\tiny\color{gray},
  columns=flexible,
  keepspaces=true,
  showstringspaces=false,
  tabsize=2,
  aboveskip=4pt,
  belowskip=4pt,
  extendedchars=true,
  literate={，}{{，}}1 {；}{{；}}1 {！}{{！}}1 {？}{{？}}1 {“}{{“}}1 {”}{{”}}1 {（}{{（}}1 {）}{{）}}1 {《}{{《}}1 {》}{{》}}1,
}

% --- Extra language definitions (no built-in listings support) ---
\lstdefinelanguage{json}{
  morekeywords={true,false,null},
  sensitive=true,
  morestring=[b]",
  comment=[l]{//},
  morecomment=[s]{/*}{*/},
}
\lstdefinelanguage{yaml}{
  morekeywords={true,false,null,yes,no,on,off},
  sensitive=false,
  morecomment=[l]{\#},
  morestring=[b]",
  morestring=[d]',
}
\lstdefinelanguage{http}{
  morekeywords={GET,POST,PUT,DELETE,PATCH,HEAD,OPTIONS,HTTP,Host,Accept,Authorization,Content-Type,Content-Length,User-Agent},
  sensitive=true,
  morecomment=[l]{\#},
  morestring=[b]",
}
\lstdefinelanguage{TypeScript}{
  morekeywords={abstract,any,as,async,await,break,case,catch,class,const,continue,debugger,declare,default,delete,do,else,enum,export,extends,false,finally,for,from,function,get,if,implements,import,in,instanceof,interface,keyof,let,module,namespace,never,new,null,number,of,package,private,protected,public,readonly,return,set,static,string,super,switch,symbol,this,throw,true,try,type,typeof,undefined,unknown,var,void,while,with,yield,Record,Array,Promise,AsyncIterable,ReadableStream,Date,Error,Object,JSON},
  sensitive=true,
  morecomment=[l]{//},
  morecomment=[s]{/*}{*/},
  morestring=[b]",
  morestring=[b]',
  morestring=[b]`{`},
}

% --- Compact spacing ---
\linespread{0.95}
\setlength{\parskip}{1pt plus 1pt minus 1pt}
\setlength{\parindent}{0pt}

% --- Table styling ---
\renewcommand{\arraystretch}{1.1}

% --- Heading styling (numbered) ---
\titleformat{\section}{\normalsize\bfseries}{\thesection\quad}{0em}{}
\titlespacing{\section}{0pt}{6pt}{2pt}
\titleformat{\subsection}{\small\bfseries}{\thesubsection\quad}{0em}{}
\titlespacing{\subsection}{0pt}{4pt}{1pt}
\titleformat{\subsubsection}{\small\bfseries}{\thesubsubsection\quad}{0em}{}
\titlespacing{\subsubsection}{0pt}{3pt}{1pt}

% --- Page style ---
\pagestyle{fancy}
\fancyhf{}
\fancyhead[L]{\small\emph{\leftmark}}
\fancyhead[R]{\small\thepage}
\renewcommand{\headrulewidth}{0.4pt}
\renewcommand{\sectionmark}[1]{}  % prevent sections from overwriting leftmark

% --- Hyperref setup ---
\hypersetup{
  colorlinks=true,
  linkcolor=blue,
  urlcolor=blue,
  citecolor=blue,
}

% --- Caption format ---
\DeclareCaptionFormat{myformat}{\small\bfseries #1#2#3}
\captionsetup{format=myformat}

\begin{document}

\title{Agent 成本与性能优化：Token 预算、限流与成本归因}
\date{}
\maketitle
\markboth{TypeScript 版 Agent 知识}{ TypeScript 版 Agent 知识}

\begin{quote}
语言策略：code（Java / Python / TypeScript）
\end{quote}



\section*{TL;DR}

\begin{enumerate}[itemsep=1pt, topsep=2pt]
  \item Agent 成本首先来自\textbf{上下文膨胀}，其次才是模型单价；先做上下文预算，再谈模型路由。
  \item 性能优化顺序：先测 P95 与 TTFT，再优化最慢的工具和最长上下文，最后才换更贵的模型。
  \item 任何优化都要绑定成本、延迟、成功率三个指标，防止“省了钱、坏了效果”。
\end{enumerate}


\section*{1. 成本公式}


\begin{lstlisting}
单任务成本 = 输入 Token × 输入单价
           + 输出 Token × 输出单价
           + 工具调用成本
           + 缓存命中节省
           + 分摊基础设施成本
\end{lstlisting}



成本归因必须带 \texttt{task\_id}、\texttt{model}、\texttt{prompt\_version}、\texttt{toolset\_version}，否则无法定位是哪个版本变贵。



\section*{2. 成本优化杠杆}

{\footnotesize
\begin{tabular}{|p{\dimexpr 0.261\columnwidth-2\tabcolsep\relax}|p{\dimexpr 0.630\columnwidth-2\tabcolsep\relax}|p{\dimexpr 0.109\columnwidth-2\tabcolsep\relax}|}
\hline
\textbf{杠杆} & \textbf{做法} & \textbf{典型收益} \\
\hline
上下文预算 & 每轮计算预算，裁剪低优先级内容 & 高 \\
\hline
模型分级 & 简单任务走小模型，复杂任务走强模型 & 高 \\
\hline
Prompt Cache & System Prompt 与工具 Schema 固定前缀 & 中高 \\
\hline
工具返回压缩 & 只保留结构化摘要 & 高 \\
\hline
轮次上限 & 降低无效循环 & 中 \\
\hline
语义缓存 & 相似问题直接复用结果 & 中，需防错 \\
\hline
批处理 & 离线任务合并调用 & 中 \\
\hline
\end{tabular}
}



\section*{3. 性能优化顺序}

\begin{enumerate}[itemsep=1pt, topsep=2pt]
  \item 建立基线：P50 / P95 / TTFT / TPOT / 单任务 Token；
  \item 找出最慢工具并设置超时与并行；
  \item 缩短上下文：裁剪历史、压缩工具返回；
  \item 流式输出降低首字延迟；
  \item 模型路由：强模型只处理复杂任务；
  \item 缓存与预取：Prompt Cache、检索预取。
\end{enumerate}


\section*{4. 成本与性能门禁}

\begin{itemize}[itemsep=1pt, topsep=2pt]
  \item 单任务成本超预算 130\% 告警；
  \item P95 超过 SLO 120\% 告警；
  \item 成本优化上线前必须过评测门禁，成功率下降即回滚；
  \item 每月输出成本归因报告：按租户、任务类型、模型、Prompt 版本、工具。
\end{itemize}


\section*{5. TypeScript 版成本追踪与限流}


\begin{lstlisting}[language=TypeScript]
interface TokenUsage {
  inputTokens: number;
  outputTokens: number;
  inputPrice: number;
  outputPrice: number;
  toolCost: number;
}

function costOf(u: TokenUsage): number {
  return u.inputTokens * u.inputPrice + u.outputTokens * u.outputPrice + u.toolCost;
}

class SlidingWindowLimiter {
  private timestamps: number[] = [];

  constructor(
    private readonly maxRequests: number,
    private readonly windowMs: number,
  ) {}

  tryAcquire(now = Date.now()): boolean {
    this.timestamps = this.timestamps.filter(t => now - t < this.windowMs);
    if (this.timestamps.length >= this.maxRequests) return false;
    this.timestamps.push(now);
    return true;
  }
}
\end{lstlisting}



\section*{6. 上线检查单}

\begin{itemize}[itemsep=1pt, topsep=2pt]
  \item [ ] 有单任务 Token 与成本预算；
  \item [ ] 有滑动窗口限流与并发上限；
  \item [ ] 有模型分级路由与降级策略；
  \item [ ] 有成本/延迟/成功率三个维度的回归对比；
  \item [ ] 缓存命中与错误率同时监控。
\end{itemize}


\section*{7. 延伸阅读}

\begin{itemize}[itemsep=1pt, topsep=2pt]
  \item \href{04-大模型网关.md}{大模型网关}
  \item \href{05-Agent可观测与追踪.md}{Agent 可观测与追踪}
  \item \href{../05-评测与质量/01-Agent评测体系.md}{Agent 评测体系}
\end{itemize}

\end{document}
